% =====================================================================
%  Introduction + Related Work (Step 3, literature-review-agent output)
%  Compiled into the master template by the section-writing-agent (Step 4).
%  Notation: \citep for parenthetical, \citet for in-text narrative.
%  cleveref is NOT available (tex_profile.json) — uses \ref / \autoref only.
% =====================================================================

\section{Introduction}
\label{sec:intro}

Large language models are now being asked the questions psychologists ask people. A fast-growing literature administers validated human personality inventories to chat-tuned LLMs, reports Big Five or Dark Triad scores, and treats the resulting numbers as evidence about an underlying machine character~\citep{miotto2023gpt3,serapio2025psychometric,jiang2023personallm,huang2023psychobench,bodroza2024personality}. These scores in turn feed into proposals to use LLMs as ``silicon samples'' that stand in for human respondents in opinion polling, behavioural experiments, and economic surveys~\citep{argyle2023one,horton2023homo,aher2023using,park2023generative}. The implicit promise is that a model that scores Conscientious on the IPIP-NEO is more conscientious, in some operational sense, than a model that does not.

This promise rests on a prior measurement question that is almost never asked: do the human instruments measure anything when they are administered to an LLM? In the human-subjects tradition, an inventory is not trusted until it passes a standard battery of psychometric checks --- internal consistency, factor structure, convergent validity, and measurement invariance~\citep{cronbach1951coefficient,kaiser1960application,vandenberg2000review}. These checks separate signal from response style, distinguish a coherent latent trait from a bag of correlated items, and tell you whether two conditions can even be compared. On LLMs these checks are rarely run; when they appear at all, they typically appear one at a time, on one inventory, on a handful of models, and are reported as numbers rather than as diagnostics~\citep{suhr2025challenging,pelet2025persistent,schaekermann2024personality}. The accumulating LLM-personality literature has therefore been writing equations of motion without checking that the ruler is rigid.

This paper runs the missing audit. We administer a 221-item battery combining four validated human inventories --- the IPIP-NEO-120 Big Five facets~\citep{johnson2014measuring}, the Short Dark Triad SD3~\citep{jones2014introducing}, the cross-cultural ZKPQ-50-CC~\citep{aluja2006cross}, and the abbreviated Eysenck EPQR-A~\citep{francis1992development} --- to 20 instruction-tuned LLMs from 8 vendor families under a Default no-persona prompt and all 16 MBTI persona prompts, with $k=5$ samples per item at temperature~$0.7$. The completed corpus contains 375{,}700 raw API calls. We then ran six standard validity checks and a persona stress test. Every check fails. The failures share a single mechanism: an alignment-induced acquiescence with a directional asymmetry --- a positive forward-minus-reverse agree-rate gap on normal-personality items (IPIP $\Delta=+0.29$) and a negative gap on dark-trait items (SD3 $\Delta=-0.33$). Acquiescence and social-desirability rejection are not two phenomena~\citep{couch1960yeasayers,crowne1960social,paulhus1991response,billiet2000modeling}; they are the same alignment-driven agree-impulse pointed in opposite directions by item valence.

The downstream cascade is mechanical (see Figure~\ref{fig:causal_chain_overview}). Item-level acquiescence inflates inter-item correlations on items it agrees with and depresses them on reverse-keyed items, so Cronbach's $\alpha$ tracks reverse-item density at Spearman $\rho=-1.0$ across the five IPIP Big Five domains. Inflated inter-domain correlations collapse the expected five-factor Big Five structure into three factors, universally across all 20 models, with Tucker's $\varphi=0.993$ and exact replication under leave-one-model-out, per-family, and per-model exploratory factor analysis. By the model level, sequential variance decomposition leaves the Model term at $0.35\%$ of total sum of squares on the Likert composite and $0.51\%$ on the binary composite --- cross-model personality comparisons are statistically meaningless at this resolution. Calibration against synthetic baselines on the same item structure makes the diagnosis explicit: LLM observed alphas hug the Random ($\alpha\!\approx\!0$) and Pure-Acquiescence ($\alpha\!\approx\!0$) baselines and stay far below a Trait+Acquiescence simulation ($\alpha>0.99$). LLMs do not look like a personality with acquiescence on top; they look like acquiescence with no personality underneath.

A persona stress test rules out the obvious rescue --- that LLMs only look uniform under a neutral instruction and that personality emerges when the model is asked to play a character. MBTI prompts move profiles by more than one standard deviation per dimension (mean Persona Separation Degree $=1.198$) and the resulting profiles are $99.7\%$ nearest-centroid recoverable (319/320 (model, persona) pairs). But cross-scale coherence on the four overlapping constructs is only $0.839$, forward/reverse pair coherence is $0.696$, and the three-factor collapse persists inside every persona condition. Persona steering replaces alignment-driven acquiescence with persona-driven compliance; it does not restore measurement validity. This dissociation between role compliance and trait variance has been hinted at in recent self-report-versus-behaviour comparisons~\citep{han2025personality}; the present audit gives it a quantitative form.

Our contributions are:
\begin{itemize}\itemsep0pt
  \item \textbf{The first multi-model, multi-instrument psychometric validation study on LLMs.} 20 models $\times$ 8 vendor families $\times$ 4 validated instruments $\times$ 17 personas $\times$ $k=5$ samples = 375{,}700 raw API calls, orders of magnitude larger than prior LLM-personality work and explicitly diagnostic rather than profiling.
  \item \textbf{A single-cause causal chain from item to model.} Item-level acquiescence $\rightarrow$ alpha tracks reverse density at $\rho=-1.0$ $\rightarrow$ EFA collapses five factors into three with Tucker's $\varphi=0.993$ $\rightarrow$ inter-model variance $<1\%$.
  \item \textbf{The directional-asymmetry framing} that unifies sycophancy~\citep{perez2023discovering,sharma2024sycophancy} and safety-driven rejection~\citep{bai2022constitutional} as two faces of one alignment artefact.
  \item \textbf{A persona stress test} that quantifies role compliance without measurement validity (PSD~$1.198$, $99.7\%$ nearest-centroid recovery, cross-scale coherence~$0.839$, Fwd/Rev coherence~$0.696$).
  \item \textbf{Synthetic-baseline calibration} that rules out a ``trait $+$ style noise'' generative model.
  \item \textbf{An item-level plasticity map} of the alignment safety floor (visible-behaviour items shift up to $0.63$ under persona prompts; safety- and morality-anchored items do not move).
  \item \textbf{A methodological prescription} for future LLM-personality work: report internal consistency, factor structure, convergent validity, and measurement invariance before interpreting any score; establish invariance before cross-condition comparisons; adopt multi-temperature and multi-seed designs.
\end{itemize}

\noindent The schematic in Figure~\ref{fig:validation_battery_schema} summarises the six-check battery and the persona stress test that organise the rest of the paper.

\section{Related Work}
\label{sec:related}

\subsection{LLM personality profiling and silicon samples}
\label{sec:related-profiling}

The dominant strand of work on LLM personality administers a human inventory to one or more models and reports the resulting profile. \citet{miotto2023gpt3} kick off this line with the NEO-PI-R and a Schwartz values survey on GPT-3. \citet{serapio2025psychometric} formalise a psychometric framework for the IPIP-NEO on Google's PaLM family and report shapable Big Five profiles. \citet{jiang2023personallm} extend the same template to PersonaLLM. \citet{huang2023psychobench} consolidate thirteen scales into the PsychoBench suite. \citet{bodroza2024personality} run temporal-stability tests across several frontier models, and \citet{lee2025trait} release the TRAIT testset built explicitly around the Big Five. Adjacent work brings cognitive-psychology probes~\citep{binz2023psychology} and trait-descriptor recovery~\citep{suh2024rediscovering} to the same population.

A parallel literature converts these profiles into operational predictions. \citet{argyle2023one} propose that LLMs conditioned on demographic backstories can simulate human survey samples; \citet{horton2023homo} extends the proposal to economic experiments; \citet{aher2023using} replicate classic behavioural studies on simulated human cohorts; \citet{park2023generative} embed persona-conditioned LLMs in a small interactive society. The silicon-sample programme inherits whatever measurement error sits in the underlying personality scores: if the inventories do not measure traits on LLMs, the simulated survey responses cannot be trusted either~\citep{tjuatja2023response,salecha2024llm}.

A small but growing critical literature has begun to push back. \citet{acerbi2024llm} and \citet{salecha2024llm} document social-desirability bias on Big Five surveys. \citet{suhr2025challenging} challenge the validity of LLM personality tests at the Web Conference. \citet{pelet2025persistent} report instability across scale, reasoning, and conversation history. \citet{hagendorff2024machine} and \citet{schaekermann2024personality} report partial reliability checks. \citet{han2025personality} expose a dissociation between LLM self-reports and downstream behaviour. \citet{xie2025aipsychobench} introduce a psychometric-difference benchmark, and~\citet{chen2025humanizing} survey the area. Common to almost all of this work is a single-instrument design, a single check (typically test-retest stability or Cronbach's $\alpha$), and a handful of models. The directional asymmetry of acquiescence between normal-personality and dark-trait inventories, the universality of the three-factor collapse under multi-instrument exploratory factor analysis, and the cross-model variance shrinkage to below one per cent of total sum of squares are, to our knowledge, not previously documented in a single audit. Our six-check battery plus persona stress test on twenty models and four instruments closes this gap.

\subsection{Response styles, acquiescence, and social desirability}
\label{sec:related-styles}

Human survey methodology has spent six decades cataloguing and controlling response styles. \citet{couch1960yeasayers} introduced acquiescence (``yea-saying'') as a personality variable; \citet{crowne1960social} developed the canonical social-desirability scale; and \citet{paulhus1991response} synthesised the field into a unified taxonomy. \citet{billiet2000modeling} model acquiescence as an explicit factor in structural-equation models of balanced item sets, and \citet{baumgartner2007modeling} document its cross-national footprint in marketing research. The standard countermeasure is reverse-keyed items: items written so that endorsement maps to the negative pole of the trait. Reverse-keyed items act both as a corrective on individual responses and as a diagnostic --- their inconsistency with forward items quantifies how much of the variance is response style rather than substance.

The diagnostic apparatus built around this idea is the same toolkit our audit imports wholesale. Cronbach's $\alpha$~\citep{cronbach1951coefficient} provides per-domain internal consistency; classical test theory~\citep{lord1968statistical} provides the underlying variance partition; exploratory factor analysis under the Kaiser eigenvalue criterion~\citep{kaiser1960application}, Horn's parallel analysis~\citep{horn1965parallel}, and varimax rotation~\citep{kaiser1958varimax} recovers latent structure; Tucker's coefficient of congruence~\citep{tucker1951method} quantifies replication; \citet{meredith1993measurement} formalised measurement invariance, and \citet{vandenberg2000review} compiled the operational checklist that organisational research now follows. Forced-choice and ipsative inventories built on item-response theory have been proposed as structural alternatives that bypass response styles entirely~\citep{brown2011forcedchoice,treaux2025forced}. Bootstrap inference~\citep{efron1979bootstrap} and Mahalanobis covariance distance~\citep{mahalanobis1936generalised} round out the statistical toolkit our audit uses.

Two limitations of the existing LLM literature stand out against this backdrop. First, the response-style toolkit has not been imported in coherent form: where reverse-keyed inconsistency is reported~\citep{acquiescence2025emnlp}, it is reported as a curiosity rather than as a quantitative diagnostic with a falsifiable structural signature. Second, the directional asymmetry between agree-pull on normal-personality items and reject-pull on dark-trait items has not been documented at the item level. We argue these are symptoms of the same alignment-driven mechanism, observable simultaneously through the same forward/reverse decomposition, and that this unification is precisely what closes the chain from item-level acquiescence to model-level homogenisation.

\subsection{Alignment, RLHF, and behavioural side effects}
\label{sec:related-alignment}

Almost every commercial LLM in our study has been trained with some combination of supervised instruction tuning~\citep{wei2022flan,ouyang2022training}, reinforcement learning from human or AI feedback~\citep{christiano2017deep,ouyang2022training,bai2022constitutional}, and direct preference optimisation~\citep{rafailov2023dpo}. The behavioural side effects of these regimens are increasingly well documented. \citet{perez2023discovering} use model-written evaluations to surface a wide spectrum of value-laden behaviours, including sycophancy and political bias. \citet{sharma2024sycophancy} characterise sycophancy as a robust failure mode that survives across multiple alignment recipes. \citet{kirk2023understanding} show that RLHF systematically reduces output diversity and homogenises generation distributions. \citet{zheng2024agreeableness} quantify agreeableness-driven sycophancy in role-playing contexts, and \citet{wang2024aligning} probe the limits of RLHF-based alignment. Survey-design probes corroborate this picture: \citet{tjuatja2023response} show that LLMs exhibit human-like response biases on survey items, and \citet{acerbi2024llm,salecha2024llm} demonstrate alignment-shaped social-desirability bias on Big Five inventories.

What this strand has not yet done is connect these side effects to classical response-style theory and show that they produce a single mechanistically traceable failure of measurement validity across the standard psychometric checks. The closest precedent is the survey-design line~\citep{tjuatja2023response,dominguez2024questioning} and the recent self-report-versus-behaviour dissociation~\citep{han2025personality}, but neither closes the loop from item-level acquiescence to model-level variance collapse. We treat alignment-induced acquiescence and its directional inversion on dark-trait items as the unified mechanism that explains (i) the $\alpha$--reverse-density rank-1 inverse relationship, (ii) the universal three-factor collapse with Tucker's $\varphi=0.993$, (iii) inter-model variance below one per cent, (iv) the item-level plasticity dichotomy between visible-behaviour items and a safety floor, and (v) persona compliance without measurement validity. Reading personality scores off chat-tuned LLMs without first auditing the instrument therefore measures the alignment layer at least as much as it measures any putative trait~\citep{serpell2024validity,laverghetta2022predicting}.
